\section{题10：不互质游戏（P16792）}

\subsection{问题重述与博弈建模}

桌面上有 $n$ 个正整数 $a_1, \dots, a_n$，两人轮流操作。每次操作选择一个数 $x$，将其替换为更小的 $y$，要求 $y < x$ 且 $\gcd(x, y) > 1$。无法操作者输。多组数据，$T \le 200$，每组 $n \le 1000$，$a_i \le 300000$，所有数据的 $n$ 之和不超过 $100000$。

这是一个公平组合游戏：双方可执行的操作完全相同，不存在随机因素，局面信息完全公开；每个桌面数字构成一个独立的子游戏，整个游戏是这些子游戏的直和。根据 Sprague–Grundy 定理，直和局面的 Grundy 数等于各子游戏 Grundy 数的异或和，先手必胜当且仅当该异或和非零。因此核心任务是为每个可能的 $a_i$（即 $1$ 到 $300000$ 的所有整数）预计算其子游戏的 Grundy 数 $G(a_i)$，其中 $G(x)$ 定义为：
\[
G(x) = \operatorname{mex}\{\, G(y) \mid 1 \le y < x,\ \gcd(x, y) > 1 \,\}.
\]
这里 $\operatorname{mex}$（minimum excluded value）表示集合中未出现的最小非负整数。边界：$G(1) = 0$（$1$ 无合法后继），对任意质数 $p$ 也有 $G(p) = 0$（不存在小于 $p$ 且与 $p$ 不互质的数）。

\subsection{朴素思路——暴力枚举为何不可行}

最直接的递推思路是：对每个 $x = 2, 3, \dots, 300000$，枚举所有 $y < x$，检查 $\gcd(x, y) > 1$ 是否成立，收集符合条件的 $G(y)$ 再求 $\operatorname{mex}$。对于 $x = 300000$，需要检查约 $3 \times 10^5$ 个候选 $y$，累计约 $\frac{1}{2} \times (3 \times 10^5)^2 \approx 4.5 \times 10^{10}$ 次 $\gcd$ 调用。即使辗转相除法每次只需 $\mathrm{O}(\log M)$ 时间，总运算量已远超 $10^{11}$ 量级，在任何评测环境下都无法通过。

可以做初步优化：不枚举所有 $y$，而是枚举 $x$ 的每个奇质因子 $p$ 的倍数。设 $x$ 的相异奇质因子为 $p_1, \dots, p_k$，合法 $y$ 的集合恰为 $\bigcup_i \{\, p_i, 2p_i, 3p_i, \dots \mid p_i \cdot t < x \,\}$。但最坏情况下 $x$ 含有多个小质因子时合法 $y$ 的数量仍可达到 $\mathrm{O}(x)$。例如 $x = 2310 = 2 \times 3 \times 5 \times 7 \times 11$，仅质因子 $3$ 就贡献了近 $770$ 个合法后继。更重要的是，我们需要的并非 $y$ 本身，而是 $G(y)$ 的集合——这一观察暗示我们可以“按质因子汇总 Grundy 值”，这正是后续位集优化的核心动机。

\subsection{偶数的封闭公式——归纳与上界证否}

偶数情形最为规整。设 $x = 2k$。所有更小的偶数 $2, 4, \dots, 2(k-1)$ 均与 $x$ 共享因子 $2$，因此全部构成合法后继。递推前几项：$G(2) = \operatorname{mex}\varnothing = 0$；$G(4) = \operatorname{mex}\{0\} = 1$；$G(6) = \operatorname{mex}\{0, 1\} = 2$；$G(8) = \operatorname{mex}\{0, 1, 2\} = 3$。归纳可知前 $k-1$ 个偶数的 Grundy 数恰好为 $0, 1, \dots, k-2$，因此偶数后继覆盖了所有 Grundy 值从 $0$ 到 $k-2$ 的连续区间。

接下来须论证 Grundy 值 $k-1$ 是否可达。若存在某个奇数 $y < 2k$ 使得 $G(y) = k-1$，则 $G(2k)$ 至少为 $k$，公式将不成立。关键观察：对任意正整数 $y$，恒有 $G(y) \le y/2 - 1$——$\operatorname{mex}$ 不可能超过集合的大小，而 $y$ 的后继 Grundy 值集合的规模至多为 $y/2 - 1$。对任意不超过 $2k-1$ 的奇数 $y$，$G(y) \le (2k-1)/2 - 1 = k - 1.5$，取整为 $k-2$；对小于 $2k$ 的偶数 $y$，$G(y) = y/2 - 1 \le (2k-2)/2 - 1 = k-2$。因此 $k-1$ 永不可达。于是得到简洁的封闭公式：
\[
G(2k) = k - 1 = \frac{x}{2} - 1.
\]
验证 $x = 12 = 2 \times 6$：公式给出 $G(12) = 5$。偶数后继 $2, 4, 6, 8, 10$ 贡献 $\{0, 1, 2, 3, 4\}$；奇数后继 $3$（$\gcd=3$）和 $9$（$\gcd=3$）贡献 $\{0, 1\}$；全体集合为 $\{0, 1, 2, 3, 4\}$，$\operatorname{mex}=5$，吻合。

\subsection{含因子 $3$ 的奇数的公式——从小数枚举中发现模式}

奇数质数的 Grundy 数为 $0$ 是平凡的。对于奇数合成数，先考察含因子 $3$ 的情况。手工计算前几项：$G(3) = 0$（质数）。$x = 9 = 3^2$：合法后继为 $3, 6$，$G(3) = 0$，$G(6) = 2$，$\operatorname{mex}\{0, 2\} = 1$。$x = 15 = 3 \times 5$：合法后继有 $3, 5, 6, 9, 10, 12$，Grundy 值为 $\{0, 0, 2, 1, 4, 5\}$，$\operatorname{mex} = 3$。$x = 21 = 3 \times 7$：后继 $3, 6, 7, 9, 12, 14, 15, 18$ 对应 Grundy $\{0, 2, 0, 1, 5, 6, 3, 8\}$，$\operatorname{mex} = 4$。$x = 27 = 3^3$：后继涵盖 $3$ 和 $2$ 的所有 $3$ 的倍数达到基本全覆盖，经过细致枚举可得 $\operatorname{mex} = 6 = \lfloor 24/4 \rfloor$。

将这些结果排成一列：$G(9) = \lfloor 6/4 \rfloor = 1$，$G(15) = \lfloor 12/4 \rfloor = 3$，$G(21) = \lfloor 18/4 \rfloor = 4$，$G(27) = \lfloor 24/4 \rfloor = 6$。模式清晰浮现：
\[
G(x) = \left\lfloor \frac{x - 3}{4} \right\rfloor, \qquad \text{对奇数 $x$ 且 $3 \mid x$}.
\]
该公式的深层逻辑是：$3$ 的倍数的 Grundy 值分布形成了缺口恰在 $\lfloor (x-3)/4 \rfloor$ 处的连续前缀。即使 $x$ 额外含有其他质因子（如 $15$ 含 $5$，$21$ 含 $7$，$45$ 含 $5$），这些额外质因子的倍数所提供的 Grundy 值全部落在上述前缀内部，永远无法精准地填补缺口。其根本原因是一条关于互质性的不动点原理：所有 Grundy 值为 $\lfloor (x-3)/4 \rfloor$ 的数必然与 $3$ 互质（否则它本该由某个 $3$ 的倍数更早取得），而既然它与 $3$ 互质，任何含有因子 $3$ 的数 $x$ 都无法将它作为合法后继。

\subsection{一般奇合成数——位集维护与按质分治}

对于最小质因子至少为 $5$ 的奇合成数（如 $25, 35, 49, 55, 77, \dots$），没有统一的封闭公式，必须回到递推框架。令 $P = \{p_1, \dots, p_k\}$ 为 $x$ 的所有互异奇质因子。由于所有 $y < x$ 的 $G(y)$ 均已计算完毕，从 $x$ 可达的 Grundy 值集合恰为“每个 $p \in P$ 的全部已处理倍数所取得的 Grundy 值的并集”。于是我们为每个奇质数 $p$ 维护一个集合 $\mathit{reach}[p]$，处理 $x$ 时求所有 $\mathit{reach}[p]$（$p \in P$）的并集的 $\operatorname{mex}$，处理完再将 $G(x)$ 纳入各 $\mathit{reach}[p]$。

对所有 $25997$ 个奇质数各维护独立集合将导致巨额开销，必须按质数大小分治。将奇质数分为“小质数”（$\le 1000$）和“大质数”（$> 1000$）两类。$M = 300000$ 内约 $168$ 个小质数，每个的倍数多达 $M/p \ge 300$ 个，适合用 \mil{std::bitset<MAX_G>} 维护——每个 \mil{bitset} 约 $19$ KB，总内存约 $3$ MB。合并时直接按位或，利用 $64$ 位机器字并行，$k$ 个 \mil{bitset} 的合并代价约为 $\mathit{MAX\_G} \cdot k / 64$ 次字操作。大质数约有 $25600$ 个，但其倍数极度稀疏（$M/p \le 300$），对每个大质数以 \mil{std::vector<pair<int,int>>} 按处理顺序记录（倍数，Grundy 值），合并时顺序扫描并筛选 \mil{multiple < x} 的条目追加到局部位集。

以 $x = 25 = 5^2$ 为例演示递推过程。处理 $x = 25$ 时，$\operatorname{spf}(25) = 5$（小质数），奇质因子仅有 $\{5\}$。此时 $\mathit{reach}[5]$ 中已有 $5$ 的倍数 $5, 10, 15, 20$ 的 Grundy 值：$G(5)=0$，$G(10)=4$，$G(15)=3$，$G(20)=9$。因此局部位集中置位的比特为 $0, 3, 4, 9$，从 $0$ 开始搜索，遇到未置位的最小值 $1$，故 $G(25) = 1$。处理完 $25$ 后，将比特 $1$ 置入 $\mathit{reach}[5]$。再处理 $x = 35 = 5 \times 7$ 时，奇质因子为 $\{5, 7\}$（$7$ 也是小质数），局部位集取 $\mathit{reach}[5] \mid \mathit{reach}[7]$，搜索 $\operatorname{mex}$ 即得 $G(35)$。整个过程对每个 $x$ 的代价正比于其奇质因子个数乘上小质数 \mil{bitset} 的长度或大质数向量的已扫描条目数。

\subsection{预处理流程与复杂度分析}

完整预处理的执行流程如下。第一步，线性筛：对 $i = 2$ 到 $M$，若 $\operatorname{spf}[i] = 0$ 则标记其为质数并加入质数列表；用每个已发现的质数 $p$ 筛去 $i \cdot p$，同时赋予 $\operatorname{spf}$ 值，当 $p > \operatorname{spf}[i]$ 时跳出内层循环以保证线性。同时将每个奇质数分配为小质数索引或大质数索引。该步骤的复杂度为 $\mathrm{O}(M \log \log M)$，对 $M = 300000$ 耗时远小于 $0.01$ 秒。

第二步，递推计算 $G(x)$。对于偶数 $x$，$G(x) = x/2 - 1$；对于奇质数 $x$，$G(x) = 0$；对于含因子 $3$ 的奇合成数，$G(x) = \lfloor (x-3)/4 \rfloor$；这三种情况各自只需 $\mathrm{O}(1)$ 时间，覆盖了约 $85\%$ 的数。对于 $\operatorname{spf} \ge 5$ 的奇合成数（约 $2 \times 10^4$ 个），每组 $\omega(x)$ 个奇质因子（$\omega(x)$ 最多约 $6$，均摊远小）需要做 \mil{bitset} 或操作或向量扫描。小质数 \mil{bitset} 按位或的单词量约为 $\mathit{MAX\_G} / 64 \approx 2300$ 字；大质数向量的扫描总条目数不超过 $\sum_{p > 1000} M/p \approx M(\ln \ln M - \ln \ln 1000) \approx 6 \times 10^5$ 个。综合运算量在 $10^7$ 级，完全可在 $0.2$ 秒内完成。

第三步，对每组查询读入 $n$ 个数，计算 $\bigoplus_i G(a_i)$。若异或和非零则输出 \mil{"Yes"}，否则输出 \mil{"No"}。$\sum n \le 100000$ 保证所有查询的累加只需线性时间。

空间复杂度：小质数 \mil{bitset} 数组约 $3$ MB；\mil{spf} 数组和 $G$ 数组各约 $1.2$ MB（\mil{int} 型）；大质数向量总条目数不超过各组互异奇质因子的出现次数之和，约在百万量级，每个条目 $8$ 字节故约 $8$ MB；总计不超过 $15$ MB。

\subsection{实现细节与注意事项}

线性筛内层循环中，乘法 \mil{i * p} 需使用 \mil{long long} 比较以防止 \mil{int} 溢出（$300000 \times 300000 \approx 9 \times 10^{10}$ 超出了 $32$ 位有符号整数的范围）。提取 $x$ 的互异奇质因子时，反复除以 $\operatorname{spf}[tmp]$ 并在 \mil{last_p} 与当前质因子相同时跳过，这样既去除了因子 $2$ 又自动去重。大质数的 \mil{reachable_large} 向量中条目按倍数递增自然有序，这是因为倍数（即各 $x$）按处理序递增；扫描时检查 \mil{entry.multiple >= x} 直接 \mil{break}，避免了无效遍历。局部位集 \mil{seen} 在每轮循环中通过 \mil{seen.reset()} 清零，清空一个长 \mil{bitset} 的代价足以忽略。$\operatorname{mex}$ 由循环 \mil{while (seen[g]) g++;} 得到，由于绝大多数情况下 $\operatorname{mex}$ 值不大，循环迭代次数极短。最终的胜负判定须牢记 SG 定理的结论：异或和非零则先手必胜；若所有 $a_i$ 的 Grundy 值异或为零，则后手有必胜策略。

\subsection{代码实现}

\inputminted{c++}{../src/p10_not_coprime_game.cpp}
